\documentclass[conference]{../../../../Conference-LaTeX-template_10-17-19/IEEEtran}
\usepackage{amsmath,amssymb,array,booktabs,tabularx}
\usepackage{url}
\usepackage{balance}
\newcolumntype{Y}{>{\raggedright\arraybackslash}X}
\title{Conditional Universals in Collective Motion: A Mechanism-Observable Program for Cross-Scale Prediction}
\author{\IEEEauthorblockN{Research Proposal}\IEEEauthorblockA{Evidence-bounded four-agent synthesis | Design-only study plan}}
\begin{document}
\maketitle
\begin{abstract}
Bird flocks, fish schools, insect swarms, bacterial suspensions, cells, and robot teams can all display extended spatiotemporal coherence. The visual resemblance of these patterns motivates the search for fundamental principles, but it does not by itself establish one universal microscopic interaction rule. This proposal develops a stronger and testable question: can a shared, conditional state description predict collective-motion observables across domains while retaining the distinct interaction channels that each domain requires? The Survey stage finds recurring roles for activity, neighbor coupling, noise, response delay, and geometry, alongside an important warning: polar order, correlation, information propagation, and causal influence are different quantities. The Idea stage selects MOSAIC, a Mechanism-Observable State Atlas and Invariance Calibration program. MOSAIC compares a hierarchical shared latent-state model against polarization-only and domain-specific baselines using blocked time, group, and domain splits. Its proposed outputs are short-horizon forecast score, regime calibration, polarization, connected correlation, correlation length, propagation lag, density heterogeneity, and vorticity or defect statistics when appropriate. The ExperimentDesign stage specifies pre-registered release gates, negative controls, ablations, and governance requirements. The result is a direct but bounded research claim: a fundamental principle of collective motion should be accepted only as a conditional predictive invariant that survives cross-domain falsification, not as a visually appealing universal law.
\end{abstract}
\begin{IEEEkeywords}
collective motion, active matter, flocking, swarm intelligence, cross-scale inference, model validation, trajectory analysis
\end{IEEEkeywords}

\section{Introduction}
Collective motion is one of the most recognizable forms of biological and engineered organization. A flock can turn coherently, a fish school can remain cohesive while changing speed, a midge swarm can display long correlations without global alignment, and a suspension of microscopic swimmers can develop coherent flow structures. The breadth of these phenomena makes collective motion an unusually productive meeting point for statistical physics, behavioral biology, fluid mechanics, control, robotics, applied mathematics, and data science. It also creates a familiar scientific danger. Similar large-scale images tempt an explanation in which every group follows the same local rule. That inference is much stronger than the evidence warrants.

The central aim of this proposal is to replace an ambiguous question--``what is the one fundamental law?''--with a falsifiable program. A useful principle should identify observables that can be measured across systems, declare the mechanisms and moderators that the principle needs, predict data that were not used to fit it, and state what outcome would reject it. The proposed answer is therefore a conditional one. Activity, coupling, noise, response delay, and geometry are candidate organizing variables, but their operational definitions and coefficients can be domain-specific. The strong claim to test is not common microscopic cognition or common hydrodynamics. It is whether a common \emph{state semantics} yields useful prediction after each system's observation channel is represented explicitly.

This distinction matters for both basic science and applications. In biology, an interaction rule that reproduces trajectories may still be a proxy for unobserved sensory, motivational, or environmental processes. In active-matter physics, a continuum theory can express symmetry and transport consequences without identifying a behavioral decision algorithm. In aerial robots, a controller must confront delay, obstacle avoidance, communication limits, and safety constraints that an idealized point-particle model omits. A cross-scale theory becomes scientifically valuable when it says which features transfer, which do not, and how to tell the difference.

The report follows the required pipeline. Section II surveys evidence and records the research gaps. Section III turns those gaps into the MOSAIC idea. Sections IV and V give a design-only benchmark and its decision rules. Sections VI--VIII delimit interpretation, identify human review requirements, and provide traceability. No trajectories are acquired, no organisms or robots are manipulated, no simulations are executed, and no observed result is claimed.

\section{Survey Evidence and Research Gap}
\subsection{From minimal models to nonequilibrium state descriptions}
Dry aligning active-matter models show that self-propulsion, local alignment, and noise can generate ordered motion, density structures, and anomalous fluctuations \cite{chate,ginelli}. These models have enduring value because they force a precise statement of a proposed mechanism. They should not, however, be mistaken for complete biological descriptions. Active systems continuously consume and dissipate energy; equilibrium constructs such as a free-energy minimization argument or detailed balance are generally unavailable \cite{ramaswamy,gompper}. Consequently, conservation laws, symmetries, advective transport, interaction range, and noise all matter to a macroscopic description.

The first survey conclusion is that a collective state must be represented by more than a single order parameter. Let the normalized velocity or heading of agent $i$ in a window be $\mathbf{u}_i$. Polarization can be written as
\begin{equation}
P=\left\|\frac{1}{N}\sum_{i=1}^{N}\mathbf{u}_i\right\|.
\end{equation}
It distinguishes a globally aligned group from a disordered one, but it does not identify the mechanism, measure directional information transfer, or determine whether correlations are long-ranged. A connected correlation function, for example,
\begin{equation}
C(r)=\left\langle\delta\mathbf{u}_i\cdot\delta\mathbf{u}_j\right\rangle_{\|\mathbf{x}_i-\mathbf{x}_j\|\in r},
\quad \delta\mathbf{u}_i=\mathbf{u}_i-\langle\mathbf{u}\rangle,
\end{equation}
captures a different property. Its length scale, a propagation lag, density heterogeneity, and rotational structure may each change without a simple change in $P$.

\subsection{Empirical interactions are not interchangeable}
The empirical literature makes this distinction unavoidable. Fish trajectory analysis inferred dynamically modulated speed and turning responses, found that directional changes can be copied from fish ahead, and reported a substantial role for three-body effects in small groups \cite{katz}. The associated modeling review warns that frequently used behavioral assumptions remain uncertain and calls for experiment-based re-examination \cite{lopez}. These papers motivate inference from trajectories, but they do not license a claim that the fitted effective force is the animal's complete causal decision process.

Bird-flock work treats turning and collective response as scientific objects in their own right \cite{cavagnaBird,cavagnaTurn}. Midge swarms give a complementary lesson: they can exhibit strong correlations despite no global collective order \cite{attanasi}. Therefore, a report of either high polarization or high correlation alone cannot settle whether a system is collectively responsive, near an ordering transition, or merely affected by a shared external driver. The Survey Agent records this as Gap G2: order, correlation, responsiveness, and causal influence must remain separate in both models and conclusions.

At microscopic scales the necessary mechanism set can differ further. Reviews of microswimmers and bacterial hydrodynamics show that low-Reynolds-number flow, propulsion geometry, boundaries, and environmental conditions can shape individual and collective behavior \cite{elgeti,lauga}. The 2020 active-matter roadmap likewise spans systems from nano- and microswimmers to animal and human scales while emphasizing that their nonequilibrium physics is broad and heterogeneous \cite{gompper}. A good cross-scale framework must leave room for those differences rather than erasing them in the name of universality.

\subsection{Engineering evidence and the gap ledger}
Engineering provides a useful stress test. Autonomous drone flocking in confined environments required explicit treatment of communication, delay, barriers, perturbations, collision avoidance, and constrained motion \cite{vasarhelyi}. Those constraints make an engineered swarm valuable for validating a controller, but success there is not direct evidence about animal sensing or bacterial hydrodynamics. The transferable lesson is methodological: a credible principle must preserve the context that changes the state.

\begin{table*}[t]
\caption{Survey Gap Ledger and Downstream Commitments}
\centering\small
\begin{tabularx}{\textwidth}{p{0.08\textwidth}Y Y Y}
\toprule
ID & Evidence-supported gap & Consequence of ignoring it & MOSAIC commitment\\
\midrule
G1 & No cross-scale map has been validated as a shared predictive state description. & Similar pictures become false evidence of a common mechanism. & Compare conditional shared and domain-specific models on held-out domains.\\
G2 & Polar order, correlation, and response propagation are distinct. & A binary ``flocking'' label conceals informative states. & Use a joint outcome vector and report each component.\\
G3 & Same-data fitting can reward visually plausible but nonpredictive models. & Inference becomes post-hoc storytelling. & Block time, group episodes, and perturbation contexts before fitting.\\
G4 & Delay and geometry are often omitted from comparisons. & Apparent universality can be a context artifact. & Encode them as explicit moderators and ablate them.\\
\bottomrule
\end{tabularx}
\end{table*}

\section{MOSAIC: The Selected Idea}
\subsection{From universal rule to conditional invariant}
The Idea Agent evaluated five directions: a single alignment-number law, taxon-specific kernel descriptions, a correlation-only classifier, a drone bridge, and MOSAIC. A single coefficient law was rejected because the survey does not support transferring one interaction strength across sensing, fluid, and control architectures. A correlation-only approach was rejected because it violates G2. The selected direction, MOSAIC--Mechanism-Observable State Atlas and Invariance Calibration--keeps a sharp scientific objective while preserving mechanism uncertainty.

For every time window $w$, MOSAIC proposes a state vector
\begin{equation}
\mathbf{z}_w=[A_w,K_w,N_w,D_w,G_w],
\end{equation}
where $A$ is activity or directional persistence, $K$ is an inferred coupling proxy with uncertainty, $N$ is residual stochasticity separated from tracking noise, $D$ is response delay, and $G$ is geometry and environmental context. The outcome vector is
\begin{equation}
\mathbf{y}_w=[P,C,\xi,\tau_{\mathrm{prop}},H,\Omega]_w,
\end{equation}
where $\xi$ is correlation length, $\tau_{\mathrm{prop}}$ is a propagation-lag estimate, $H$ captures density heterogeneity, and $\Omega$ represents vorticity or defect statistics only for modalities where it is meaningful.

These symbols do not assert that every modality measures the same physical quantity in the same units. Instead, each domain has an explicit observation map $g_d$:
\begin{equation}
\mathbf{y}_{w,d}=g_d(\mathbf{z}_{w,d},\mathbf{q}_{w,d})+\boldsymbol{\epsilon}_{w,d},
\end{equation}
where $\mathbf{q}$ records data quality, sampling rate, field of view, and known context. A shared latent transition model is tested only after $g_d$ and $\mathbf{q}$ are specified. This is the core safety against a spurious cross-species ``collapse.''

\subsection{Scientific claims and failure conditions}
MOSAIC makes three claims that can fail. First, the joint state should predict collective regime and short-horizon trajectories more accurately than polarization alone. Second, a hierarchical conditional model should be practically non-inferior to a domain-specific model on at least one held-out domain only if delay and geometry are retained. Third, if a shared model fails its predictive release gate or needs incompatible state semantics across domains, the cross-scale invariant claim is rejected. These failure conditions are part of the result, not errors to conceal.

The proposed principle is therefore: \emph{collective coherence is organized by conditional balances among persistent activity, coupling, fluctuations, latency, and geometry, and a proposed balance earns fundamental status only through predictive invariance under context-aware validation.} This is more direct than an all-purpose disclaimer, yet it is falsifiable and does not claim that every agent processes the same information.

\section{Experiment Design: A Design-Only Cross-Scale Benchmark}
\subsection{Cohorts, units, and acquisition boundary}
The initial study is a retrospective benchmark using lawfully reusable public trajectory or velocity-field data. Its planned domains are three-dimensional bird or fish trajectories, two-dimensional fish schools, bacterial or microswimmer velocity fields, and robot-swarm logs or a future validated simulator. A domain may be substituted only with a documented license, metadata card, preprocessing record, and measurement-quality assessment. The experimental unit is a fixed-duration observation window containing agent trajectories or a velocity field, not an individual point in a shuffled trace.

The design is deliberately digital in its initial phase. It neither requests a new animal experiment nor proposes autonomous drone operation. Any later acquisition involving animals requires protocol, welfare, permitting, and location-risk review. Bacterial or microswimmer acquisition requires trained laboratory supervision and applicable biological-safety procedures. Drone work requires human-supervised hardware, airspace, and safety review. None of those steps has been taken here.

\subsection{Measurement and inference plan}
For every window, preprocessing records coordinate system, sample rate, missingness, smoothing policy, group size, field of view, and boundary geometry. Headings are estimated only when temporal resolution supports a stable derivative; otherwise the analysis reports a coarser velocity-field observable. Candidate neighbor kernels include metric distance, fixed-count neighbors, visually constrained neighbors where visibility can be reconstructed, and field-mediated couplings for applicable fluid data. The output is a posterior or uncertainty interval over kernel summaries rather than a single asserted true graph.

Activity is represented by speed and directional-persistence summaries. Coupling is represented by a fitted response gain to specified neighbor statistics. Noise is a conditional residual after measurement uncertainty is modeled. Delay is a lag between a local heading or velocity perturbation proxy and the best-supported neighbor response. Geometry includes obstacle proximity, confinement, curvature where relevant, local density, and environmental-flow or communication information. Dimensionless ratios such as an alignment-to-noise ratio or delay-to-turning-time ratio are permitted only as preregistered, data-specific normalizations. They cannot be selected after inspecting whether they create an attractive collapse.

\begin{table}[t]
\caption{Model Family and Test Purpose}
\centering\small
\begin{tabular}{p{0.10\columnwidth}p{0.24\columnwidth}p{0.38\columnwidth}}
\toprule
Model & Inputs & Test purpose\\
\midrule
M0 & Domain-specific states and kernels & Strong descriptive baseline; no transfer assumption.\\
M1 & Polarization and density & Tests whether a single global descriptor is sufficient.\\
M2 & MOSAIC latent state $[A,K,N,D,G]$ plus domain observation maps & Tests conditional cross-domain transfer.\\
M3 & M2 with one mechanism removed at a time & Tests whether delay, geometry, coupling uncertainty, and multi-observable outputs are material.\\
\bottomrule
\end{tabular}
\end{table}

The primary prediction task is short-horizon heading or velocity forecasting, scored by a proper probabilistic log score or calibrated angular-error distribution appropriate to the data. The secondary task predicts the joint collective regime vector $\mathbf{y}$. The report will show absolute and relative performance, calibration curves, and all domain-specific results. It will not pool domains into a single favorable average before revealing failures.

\subsection{Splits, controls, and analyses}
Randomly shuffling individual points makes temporal autocorrelation leak across train and test sets. MOSAIC instead blocks contiguous time windows, full group episodes, and, where labels exist, all examples of the same perturbation context. A leave-one-domain-out protocol evaluates transfer. Hyperparameters, normalizations, and candidate kernels are selected only within development domains. The held-out domain is opened once after the analysis plan is frozen.

Three negative controls are essential. A temporal-shuffle control breaks response order while retaining marginal heading distributions. A group-membership shuffle tests whether apparent interaction is attributable to a shared external drive or preprocessing artifact. A heading-shuffle control tests whether the algorithm can manufacture a convincing kernel from no directional relation. Synthetic truth-known sequences are allowed to assess recovery of a known model, but they are never used as evidence that a biological mechanism has been found.

Robustness analysis repeats the pipeline over tracking-quality thresholds, smoothing windows, field-of-view masks, missing-agent patterns, and alternate neighbor definitions. The study will report sensitivity to finite group size and derivative estimation. A domain with insufficient temporal or spatial resolution is marked uninformative for the relevant endpoint; forcing it into a pooled result is prohibited.

\subsection{Identifiability and Cross-Domain Measurement Invariance}
A predictive comparison can still be scientifically weak if two different mechanisms generate the same measured trajectories. This is the identifiability problem. An animal may change heading because of a neighbor's visible turn, a shared predator cue, an environmental flow, an unobserved goal, or the coupled effect of several neighbors. A bacterium can align through steric contact, near-field hydrodynamics, chemical gradients, confinement, or a combination whose apparent signature changes with density. A robot may copy a neighbor because a controller commands it, because a shared navigation waypoint changes, or because a collision-avoidance layer overrides its flocking term. The state vector $[A,K,N,D,G]$ does not remove these ambiguities. Its purpose is to make them visible as alternative measurement and inference models.

For this reason MOSAIC separates three levels of statement. The first is a descriptive statement: a candidate kernel and state-space model reproduce a distribution of observed headings or velocity fields. The second is a predictive statement: the same model predicts a blocked future window or a new group episode better than a baseline. The third is a mechanistic statement: an intervention that selectively changes a proposed channel produces the predicted change while important alternatives are controlled. The initial benchmark is authorized only for the first two levels. Its positive outcome may establish conditional predictive transfer. It may not promote an effective coupling parameter to a biological sensory rule, a hydrodynamic force law, or a cognitive mechanism.

The design uses an identifiability ladder. At Level 0, data quality is insufficient to estimate the endpoint and the domain is marked uninformative. At Level 1, a descriptive fit is available but multiple kernels are observationally equivalent. At Level 2, blocked forecasts distinguish candidate models within domain. At Level 3, a held-out domain retains the specified state semantics and meets the non-inferiority release gate. Level 4, which lies outside this proposal, would require a controlled and ethically approved intervention that perturbs one contextual variable while preserving the others as far as practical. Every paper, plot, and model card should state the highest supported level. This prevents a strong-looking fit from being narrated as an intervention result.

Measurement invariance is the corresponding cross-domain problem. Before interpreting a common latent state, the analysis tests whether an increase in a latent coordinate has a comparable relationship to observable outcomes after the domain-specific map is applied. Full equality of coefficients is neither expected nor required. The minimum requirement is weaker and testable: a monotonic or otherwise preregistered relationship must survive the stated calibration, and residual errors must not be systematically explained by an omitted domain label. If a latent coordinate improves one domain only because it encodes sampling rate, coordinate smoothing, or field of view, it is not an invariant mechanism proxy.

A staged calibration procedure is proposed. Stage A harmonizes units and uncertainty records without fitting a shared model. Stage B fits each domain separately and examines whether the proposed outcomes are observable with adequate reliability. Stage C fits the hierarchical model with domain-specific observation maps. Stage D performs leave-one-domain-out transfer. Stage E runs ablations that remove delay, geometry, coupling uncertainty, and all but polarization. A shared claim is released only at Stage D and only if Stage E confirms that the named variables add predictive information. The stages are deliberately ordered: forced pooling before basic observability would make a shared result impossible to interpret.

The proposal also distinguishes invariance of \emph{state} from invariance of \emph{control}. A bird flock and a robot team could occupy comparable high-polarization, short-delay state regions while reaching them through entirely different controllers. Such a result would be valuable for prediction and design, but it would not tell an engineer to copy animal behavior or a biologist to interpret a robot controller as a model of cognition. Conversely, a lack of shared state semantics would not erase useful domain-specific laws. The fundamental question is therefore plural in practice: which conditional relations remain predictive, over which range, at which identifiability level, and with what observation model?

This framework gives negative results a constructive role. If a shared latent coordinate is not invariant, investigators can inspect whether its failure tracks medium type, sensing range, confinement, group size, or data quality. The next iteration may split a state, add a justified moderator, or narrow the target domain; it may not redefine the endpoint after the fact until a transfer passes. That discipline turns heterogeneity from a nuisance into evidence about the boundary of a proposed principle.
\section{Validation and Decision Rules}
\subsection{Release gate}
The primary comparison is the held-out predictive-score difference between M2 and M0. Before data are inspected, the investigators specify a practical non-inferiority margin $\Delta_d$ for each outcome scale and an uncertainty procedure. Define
\begin{equation}
\Delta S_d=S(\mathrm{M2};d)-S(\mathrm{M0};d),
\end{equation}
with larger $S$ meaning better calibrated prediction. Conditional transfer is supported for domain $d$ only if the lower uncertainty bound for $\Delta S_d$ exceeds $-\Delta_d$, M2 improves over M1, and no required ablation renders the claimed shared state superfluous. Each condition has a scientific role: M0 prevents false confidence against a weak baseline, M1 prevents unnecessary complexity, and M3 distinguishes a real mechanism balance from a decorative feature list.

The decision has four labels. \emph{Conditional-transfer supported} means all release gates pass for a defined domain and observable set. \emph{Domain-limited} means M2 works within a domain but not in transfer. \emph{Underdetermined} means data quality or controls prevent discrimination. \emph{Conditional-transfer not supported} means M2 misses the predictive margin or fails a critical control. None of these labels states that a microscopic causal rule has been directly observed.

\subsection{What observations would mean}
If M2 transfers only when delay and geometry are included, the outcome supports a contextual rather than bare alignment principle. If M1 performs comparably, the additional mechanism proxies have not shown incremental predictive value. If negative controls also predict well, the data pipeline has insufficient discriminative validity. If M0 decisively wins in every held-out domain, the honest conclusion is that the proposed shared state semantics did not transfer. That negative result would directly narrow the scope of ``fundamental principles'' and guide the next Survey iteration toward domain-specific mechanisms.

\section{Limitations, Risks, and Human Review}
The proposal makes no claim that data-driven effective interactions equal perception, intention, molecular signaling, or hydrodynamic force. A trajectory can contain unmeasured common drivers, selection effects, tracking error, and state-dependent behavior. A hydrodynamic field and an animal's heading trace have different measurement operators; a successful latent alignment must be inspected for measurement invariance rather than accepted because a plot looks collapsed. Bacteria, microswimmers, animals, and robots also differ in energy use, sensing, embodiment, and objective functions. These differences are central moderators, not inconvenient noise.

Several review decisions require qualified humans. Curators must confirm dataset licenses, permissions, and metadata. Behavioral experts must assess whether tracking conditions and inferred kernels are biologically plausible. Fluid and active-matter specialists must determine whether field resolution supports vorticity or defect inference. Roboticists must approve any later platform validation. The frozen bibliography was cross-validated at metadata level through OpenAlex and AnySearch; direct publisher/DOI inspection remains pending because the requested in-app browser could not start under the current Windows ACL. The evidence and proposal should be updated after that review rather than quietly treating metadata as full-text verification.

\section{Conclusion}
Collective motion has no evidence-backed single law that can be read from a flock's shape alone. It does have a promising scientific structure: persistent activity, neighbor coupling, fluctuations, delay, and geometry recur as conditional organizing variables, while collective order, correlation, propagation, and prediction provide complementary measurements. MOSAIC turns that structure into a rigorous test. It asks whether a context-aware shared state model can survive held-out cross-domain prediction against domain-specific alternatives. Passing would support a conditional predictive invariant; failure would refute the proposed transfer rather than dilute the question. This is the appropriate standard for calling a principle fundamental in a multidisciplinary, nonequilibrium field.

\appendices
\section{Minimum Data and Provenance Schema}
Every analysis window should retain an immutable identifier; source dataset version and license; domain and species/platform label; coordinate frame; sampling rate; time span; number of tracked agents or field cells; missingness; smoothing and derivative method; field-of-view mask; group episode identifier; geometry and obstacle metadata; environmental or communication context; measurement uncertainty; candidate-kernel family; train/validation/test assignment; and code or workflow version. No pooled comparison should be accepted when these records are absent. This schema separates a scientific claim from a convenient but unreproducible visualization.

\section{Pre-Registered Analysis Registry}
The conditional-universality question can be compromised even when every individual model is technically competent. The principal risk is analytic flexibility: a researcher can alter the time window, tracking filter, neighbor definition, normalization, target variable, or transfer split until unrelated systems appear to share a curve. MOSAIC therefore requires a compact pre-registration registry before the final test sets are opened. The registry is not a claim that all future scientific judgment can be automated; it is a record of the choices most capable of manufacturing an apparent invariant.

First, the registry names each domain, data version, license, coordinate convention, sampling interval, spatial coverage, and inclusion or exclusion rule. It then commits to the outcome families. A trajectory domain may support angular forecasts, polarization, correlation, and response lags, whereas a velocity-field domain may support flow forecast, correlation, spectral, and vorticity quantities. An endpoint that cannot be measured at adequate resolution is declared unavailable rather than filled by a proxy after model fitting. This distinction prevents the shared latent state from being evaluated on a different outcome in every domain while retaining the same favorable label.

Second, the registry defines a finite candidate set of interaction kernels. For discrete agents, the set may include metric, fixed-count, and visibility-constrained neighborhoods. For fields or microswimmers, it may include local gradients, finite-range alignment, or explicitly stated flow-mediated summaries. The procedure evaluates each candidate within development data and reports the full comparison. It may not add a new kernel after observing a failed transfer result unless a new Survey and Idea iteration records the reason. The same rule applies to the transformation of raw coordinates into $A$, $K$, $N$, $D$, and $G$: every threshold, temporal lag range, and normalization has a declared provenance.

Third, the registry commits to the split hierarchy. Individuals are never independently randomized across training and testing when they remain in one group episode. The study blocks time, then group episode, then domain, and keeps perturbation contexts intact. It records the random seed for any within-development resampling and a single sealed held-out evaluation set. The model score is accompanied by predictive calibration, interval coverage, and an error decomposition by density, geometry, group size, and tracking quality. A pooled score is supplementary; it can never overwrite a documented domain failure.

Finally, the registry specifies the interpretation language in advance. Passing M2 against M0 within the practical margin permits only the phrase \emph{conditional-transfer supported for the stated domains, observables, and contexts}. It does not permit ``universal mechanism discovered.'' Failing the margin requires the phrase \emph{conditional-transfer not supported under this measurement model}. It does not prove that no shared physical principle exists. These semantic release rules protect the conclusion from both overclaiming and an evasive reinterpretation of a negative result.

\section{Evidence-to-Claim Traceability}
The authoring stage preserves a claim-to-source record instead of allowing a general active-matter review to function as a citation for every conclusion. The entries below describe the scope of support; none represents a claim that the cited papers directly establish the entire MOSAIC synthesis.

\noindent\textbf{T1: Candidate organizing variables.} The proposal's use of persistent activity, local coupling, fluctuations, and symmetry-aware transport as candidate state variables draws on dry-aligning and active-matter reviews \cite{chate,ramaswamy,ginelli,gompper}. Those sources motivate a conditional theoretical vocabulary. They do not demonstrate a complete sensory or behavioral model for any animal group, nor do they establish numerical coefficient equality across media.

\noindent\textbf{T2: Multi-observable state measurement.} The requirement to separate global order, correlation, and response propagation is grounded in bird-flock, turning, midge, and fish studies \cite{cavagnaBird,cavagnaTurn,attanasi,katz}. In particular, correlation without global order makes a polarization-only classification insufficient. This support is methodological: it calls for joint reporting of endpoints and does not turn a correlation estimate into proof of direct causal influence.

\noindent\textbf{T3: Inference and predictive testing.} The use of alternative neighbor kernels, held-out prediction, and caution about untested behavioral assumptions follows the fish modeling and interaction-inference literature \cite{lopez,katz}. The evidence supports the need to infer and test effective descriptions from trajectories. It does not ensure that even a held-out winner uniquely identifies the true interaction graph or an individual's decision algorithm.

\noindent\textbf{T4: Context moderators.} Microswimmer and bacterial reviews motivate retaining medium, boundary, and hydrodynamic considerations, while drone-swarm work motivates retaining delay, obstacles, and communication constraints \cite{elgeti,lauga,vasarhelyi}. MOSAIC treats these terms as moderators precisely because their meanings need not be identical in every domain. This evidence rules out suppressing context in the name of a visually simple universal law.

\section{Protocol Deviation and Null-Result Policy}
A future execution must distinguish a planned design revision from a post-hoc rescue. A protocol deviation is any change after development data are examined to the outcome definition, data-quality threshold, candidate kernel set, normalization, split assignment, practical margin, or interpretation label. Each deviation requires a dated rationale, the affected domains, an assessment of whether it could favor M2, and a rerun of all affected baselines. The original analysis remains available as a sensitivity result. A deviation that introduces a new shared variable, a new domain, or an intervention claim returns the workflow to Survey and Idea review rather than silently modifying the final test.

Null and mixed outcomes are first-class deliverables. If M0 wins in a held-out domain, the report identifies whether the failure arises in observation reliability, state calibration, forecast horizon, geometry, delay, or a specific regime. If M2 is non-inferior but M3 shows that a named variable is dispensable, the conclusion is narrowed to the smaller state. If a negative control retains predictive performance, the correct label is underdetermined, not conditional transfer. These rules ensure that a failure to find a shared invariant produces a usable map of its boundary conditions instead of an uninformative claim that collective motion is ``too complex.''

\section{Proposal Boundary}
This document is a research proposal with status \texttt{PROPOSAL\_NO\_OBSERVED\_RESULTS}. It does not report a new experiment, simulation, biological measurement, drone deployment, or discovered collective-motion law. Its deliverable is the falsifiable benchmark, its release gate, and the conditional interpretation rules described above.

\balance
\begin{thebibliography}{00}
\bibitem{chate} H. Chat\'{e}, ``Dry Aligning Dilute Active Matter,'' \emph{Annual Review of Condensed Matter Physics}, vol. 11, pp. 189--212, 2020, doi: 10.1146/annurev-conmatphys-031119-050752.
\bibitem{ramaswamy} S. Ramaswamy, ``Active matter,'' \emph{Journal of Statistical Mechanics: Theory and Experiment}, vol. 2017, p. 054002, 2017, doi: 10.1088/1742-5468/aa6bc5.
\bibitem{elgeti} J. Elgeti, R. G. Winkler, and G. Gompper, ``Physics of microswimmers--single particle motion and collective behavior: a review,'' \emph{Reports on Progress in Physics}, vol. 78, p. 056601, 2015, doi: 10.1088/0034-4885/78/5/056601.
\bibitem{cavagnaBird} A. Cavagna and I. Giardina, ``Bird Flocks as Condensed Matter,'' \emph{Annual Review of Condensed Matter Physics}, vol. 5, pp. 183--207, 2014, doi: 10.1146/annurev-conmatphys-031113-133834.
\bibitem{lopez} U. Lopez, J. Gautrais, I. D. Couzin, and G. Th\'{e}raulaz, ``From behavioural analyses to models of collective motion in fish schools,'' \emph{Interface Focus}, vol. 2, pp. 693--707, 2012, doi: 10.1098/rsfs.2012.0033.
\bibitem{ginelli} F. Ginelli, ``The Physics of the Vicsek model,'' \emph{European Physical Journal Special Topics}, vol. 225, pp. 2099--2117, 2016, doi: 10.1140/epjst/e2016-60066-8.
\bibitem{cavagnaTurn} A. Cavagna \emph{et al.}, ``Flocking and Turning: a New Model for Self-organized Collective Motion,'' \emph{Journal of Statistical Physics}, vol. 158, pp. 601--627, 2014, doi: 10.1007/s10955-014-1119-3.
\bibitem{attanasi} A. Attanasi \emph{et al.}, ``Collective Behaviour without Collective Order in Wild Swarms of Midges,'' \emph{PLoS Computational Biology}, vol. 10, p. e1003697, 2014, doi: 10.1371/journal.pcbi.1003697.
\bibitem{katz} Y. Katz, K. Tunstr\"om, C. C. Ioannou, C. Huepe, and I. D. Couzin, ``Inferring the structure and dynamics of interactions in schooling fish,'' \emph{Proceedings of the National Academy of Sciences}, vol. 108, pp. 18720--18725, 2011, doi: 10.1073/pnas.1107583108.
\bibitem{vasarhelyi} G. Vas\'{a}rhelyi \emph{et al.}, ``Optimized flocking of autonomous drones in confined environments,'' \emph{Science Robotics}, vol. 3, 2018, doi: 10.1126/scirobotics.aat3536.
\bibitem{gompper} G. Gompper \emph{et al.}, ``The 2020 motile active matter roadmap,'' \emph{Journal of Physics: Condensed Matter}, vol. 32, p. 193001, 2020, doi: 10.1088/1361-648x/ab6348.
\bibitem{lauga} E. Lauga, ``Bacterial Hydrodynamics,'' \emph{Annual Review of Fluid Mechanics}, vol. 48, pp. 105--130, 2016, doi: 10.1146/annurev-fluid-122414-034606.
\end{thebibliography}
\end{document}
